\chapter{Streuung und Formfaktoren}

\section{Ziel dieses Kapitels}

In diesem Kapitel geht es darum, wie man aus Streuexperimenten Informationen über
Kerne und Teilchen gewinnt.

Die zentrale Idee ist:
\begin{center}
\emph{Man schießt bekannte Teilchen auf ein Target und rekonstruiert aus der Ablenkung
die Struktur des Targets.}
\end{center}

\begin{notebox}
Dieses Kapitel behandelt die Stichwörter:
\begin{itemize}
    \item Rutherfordstreuung,
    \item Streupotentiale,
    \item Wirkungsquerschnitt,
    \item Impulsübertrag,
    \item Mottstreuung,
    \item Rutherfordradius,
    \item Strukturauflösung von Kernen und Teilchen,
    \item elektromagnetischer Radius,
    \item Formfaktor.
\end{itemize}
\end{notebox}

\section{Grundidee eines Streuexperiments}

Bei einem Streuexperiment lässt man ein Projektil auf ein Target treffen.
Das Projektil kann zum Beispiel ein Alpha-Teilchen, ein Elektron oder ein Proton sein.
Das Target kann ein Atomkern, ein Proton oder ein anderes Teilchen sein.

\begin{definitionbox}
Ein \emph{Streuexperiment} ist ein Experiment, bei dem ein einfallendes Teilchen an
einem Target abgelenkt wird.

Aus der Winkelverteilung und Energieverteilung der gestreuten Teilchen schließt man
auf die Wechselwirkung und die innere Struktur des Targets.
\end{definitionbox}

\begin{conceptbox}
Streuung ist in der Kern- und Teilchenphysik eines der wichtigsten Werkzeuge.

Der Grund ist:
Sehr kleine Objekte kann man nicht direkt mit Lichtmikroskopen abbilden.
Stattdessen verwendet man Teilchen mit kleiner de-Broglie-Wellenlänge als Sonden.
\end{conceptbox}

\begin{examtipbox}
Eine typische Prüfungsfrage ist:
\begin{center}
\emph{Warum kann man mit höherer Teilchenenergie kleinere Strukturen auflösen?}
\end{center}
Die Antwort ist:
Höhere Energie bedeutet größeren Impuls.
Größerer Impuls bedeutet kleinere de-Broglie-Wellenlänge.
Damit steigt die Strukturauflösung.
\end{examtipbox}

\section{Streupotentiale}

Die Ablenkung eines Projektils hängt vom Potential ab, in dem es sich bewegt.

\begin{definitionbox}
Ein \emph{Streupotential} beschreibt die Wechselwirkung zwischen Projektil und Target
als Funktion des Abstands $r$.
\end{definitionbox}

\begin{conceptbox}
In der Kernphysik treten besonders zwei Potentialtypen auf:

\begin{itemize}
    \item das Coulomb-Potential,
    \item das kurzreichweitige Kernpotential.
\end{itemize}

Das Coulomb-Potential ist langreichweitig.
Das Kernpotential ist kurzreichweitig und wirkt nur auf Abständen der Größenordnung
\[
    r \sim \SI{1}{fm}.
\]
\end{conceptbox}

\begin{formulabox}
Für zwei Punktladungen mit Ladungen
\[
    Z_1e
    \qquad\text{und}\qquad
    Z_2e
\]
ist das Coulomb-Potential
\[
    V_{\mathrm{C}}(r)
    =
    \frac{1}{4\pi\varepsilon_0}
    \frac{Z_1Z_2e^2}{r}.
\]
\end{formulabox}

\begin{conceptbox}
Bei großen Abständen dominiert die Coulomb-Wechselwirkung.

Bei sehr kleinen Abständen kann zusätzlich die starke Wechselwirkung relevant werden.
Dann weicht die gemessene Streuung von der reinen Coulomb- beziehungsweise
Rutherfordstreuung ab.
\end{conceptbox}

\section{Wirkungsquerschnitt}

In Streuexperimenten misst man nicht einzelne Bahnen, sondern Wahrscheinlichkeiten.
Diese werden durch Wirkungsquerschnitte beschrieben.

\begin{definitionbox}
Der \emph{Wirkungsquerschnitt} $\sigma$ ist ein Maß für die Wahrscheinlichkeit,
dass ein bestimmter Streu- oder Reaktionsprozess stattfindet.

Er hat die Dimension einer Fläche.
\end{definitionbox}

\begin{definitionbox}
Der \emph{differentielle Wirkungsquerschnitt}
\[
    \frac{\dd\sigma}{\dd\Omega}
\]
gibt an, wie stark in ein Raumwinkelelement $\dd\Omega$ gestreut wird.
\end{definitionbox}

\begin{formulabox}
Das Raumwinkelelement in Kugelkoordinaten ist
\[
    \dd\Omega = \sin\theta\,\dd\theta\,\dd\varphi.
\]
Der totale Wirkungsquerschnitt ergibt sich durch Integration:
\[
    \sigma
    =
    \int \frac{\dd\sigma}{\dd\Omega}\,\dd\Omega.
\]
\end{formulabox}

\begin{conceptbox}
Ein großer Wirkungsquerschnitt bedeutet:
Der Prozess ist wahrscheinlich.

Ein kleiner Wirkungsquerschnitt bedeutet:
Der Prozess ist selten.

Der Wirkungsquerschnitt ist aber keine geometrische Fläche im einfachen Sinn.
Er ist eine effektive Fläche, die die Stärke eines Prozesses beschreibt.
\end{conceptbox}

\begin{examtipbox}
Wenn in einer Aufgabe nach einer Winkelverteilung gefragt ist, ist meistens
\[
    \frac{\dd\sigma}{\dd\Omega}
\]
gemeint.

Wenn nach der Gesamtwahrscheinlichkeit oder Gesamtfläche gefragt ist, muss man über
den Raumwinkel integrieren:
\[
    \sigma
    =
    \int \frac{\dd\sigma}{\dd\Omega}\,\dd\Omega.
\]
\end{examtipbox}

\section{Rutherfordstreuung}

Die Rutherfordstreuung beschreibt die elastische Streuung geladener Teilchen an
einem Coulomb-Potential.

Historisch wurde sie wichtig, weil sie zeigte, dass die positive Ladung eines Atoms
in einem sehr kleinen Kern konzentriert ist.

\begin{definitionbox}
\emph{Rutherfordstreuung} ist die elastische Streuung eines geladenen Projektils an
einem schweren geladenen Target, wobei die Wechselwirkung durch das Coulomb-Potential
beschrieben wird.
\end{definitionbox}

\begin{conceptbox}
Für Alpha-Teilchen an einem schweren Kern gilt näherungsweise:
\begin{itemize}
    \item Das Target ist viel schwerer als das Projektil.
    \item Die Streuung ist elastisch.
    \item Die Wechselwirkung ist Coulomb-Abstoßung.
    \item Der Kern wird zunächst als punktförmig behandelt.
\end{itemize}
\end{conceptbox}

\begin{formulabox}
Für die Rutherfordstreuung gilt näherungsweise:
\[
    \frac{\dd\sigma}{\dd\Omega}
    =
    \left(
        \frac{1}{4\pi\varepsilon_0}
        \frac{Z_1Z_2e^2}{4E}
    \right)^2
    \frac{1}{\sin^4(\theta/2)}.
\]
Dabei ist:
\begin{itemize}
    \item $Z_1e$ die Ladung des Projektils,
    \item $Z_2e$ die Ladung des Targets,
    \item $E$ die kinetische Energie im Schwerpunktsystem beziehungsweise näherungsweise im Laborsystem bei schwerem Target,
    \item $\theta$ der Streuwinkel.
\end{itemize}
\end{formulabox}

\begin{conceptbox}
Die wichtigste Winkelabhängigkeit ist:
\[
    \frac{\dd\sigma}{\dd\Omega}
    \propto
    \frac{1}{\sin^4(\theta/2)}.
\]

Für kleine Winkel gilt:
\[
    \sin(\theta/2)\approx \theta/2.
\]
Damit wird der Wirkungsquerschnitt bei kleinen Winkeln sehr groß.
Die Rutherfordstreuung ist also stark vorwärtsgerichtet.
\end{conceptbox}

\begin{examtipbox}
Typische qualitative Aussage:
\[
    \theta \to 0
    \quad\Rightarrow\quad
    \frac{\dd\sigma}{\dd\Omega}\to \infty.
\]
Das bedeutet nicht, dass unendlich viele Teilchen gestreut werden.
Es zeigt, dass die Coulomb-Wechselwirkung langreichweitig ist und kleine Ablenkungen
sehr häufig sind.
\end{examtipbox}

\section{Rutherfordstreuung und Kernradius}

Die Rutherfordform gilt nur solange, wie das Projektil den Kern nicht als ausgedehntes
Objekt erkennt.

\begin{conceptbox}
Bei niedrigen Energien oder großen Abständen folgt die Streuung der Coulomb-Kurve.

Wenn die Projektilenergie groß genug wird, kann das Projektil näher an den Kern
herankommen.
Dann können Abweichungen von der Rutherfordstreuung auftreten.

Solche Abweichungen sind ein experimenteller Hinweis auf:
\begin{itemize}
    \item die endliche Ausdehnung des Kerns,
    \item die Ladungsverteilung im Kern,
    \item zusätzliche starke Wechselwirkung bei kleinen Abständen.
\end{itemize}
\end{conceptbox}

\begin{definitionbox}
Der \emph{Rutherfordradius} ist ein Radiusbegriff, der aus der kleinsten Annäherung
eines geladenen Projektils an den Kern in einem Coulomb-Streuexperiment motiviert ist.
\end{definitionbox}

\begin{formulabox}
Für eine zentrale Annäherung kann man die kleinste Entfernung grob aus Energieerhaltung
abschätzen:
\[
    E
    =
    \frac{1}{4\pi\varepsilon_0}
    \frac{Z_1Z_2e^2}{r_{\min}}.
\]
Daraus folgt:
\[
    r_{\min}
    =
    \frac{1}{4\pi\varepsilon_0}
    \frac{Z_1Z_2e^2}{E}.
\]
\end{formulabox}

\begin{conceptbox}
Wenn
\[
    r_{\min}
\]
größer als der Kernradius ist, sieht das Projektil hauptsächlich das Coulombfeld.

Wenn
\[
    r_{\min}
\]
in die Größenordnung des Kernradius kommt, kann die gemessene Streuung von der
Rutherfordform abweichen.
\end{conceptbox}

\section{Impulsübertrag}

Bei jeder Streuung ändert sich der Impuls des Projektils.
Diese Änderung nennt man Impulsübertrag.

\begin{definitionbox}
Der \emph{Impulsübertrag} $\vec q$ ist die Differenz zwischen Endimpuls und Anfangsimpuls
des gestreuten Teilchens:
\[
    \vec q = \vec p_{\mathrm{f}}-\vec p_{\mathrm{i}}.
\]
\end{definitionbox}

\begin{formulabox}
Für elastische Streuung mit
\[
    |\vec p_{\mathrm{i}}|=|\vec p_{\mathrm{f}}|=p
\]
gilt:
\[
    q = |\vec q|
      =
    2p\sin\left(\frac{\theta}{2}\right).
\]
\end{formulabox}

\begin{conceptbox}
Der Impulsübertrag ist die zentrale Größe für Strukturuntersuchungen.

Kleine Streuwinkel bedeuten kleinen Impulsübertrag:
\[
    \theta \approx 0
    \quad\Rightarrow\quad
    q\approx 0.
\]

Große Streuwinkel bedeuten großen Impulsübertrag:
\[
    \theta \text{ groß}
    \quad\Rightarrow\quad
    q \text{ groß}.
\]
\end{conceptbox}

\begin{derivationbox}
Für elastische Streuung bilden Anfangsimpuls $\vec p_{\mathrm{i}}$ und Endimpuls
$\vec p_{\mathrm{f}}$ einen Winkel $\theta$.
Da beide denselben Betrag $p$ haben, gilt geometrisch:
\[
    q^2
    =
    |\vec p_{\mathrm{f}}-\vec p_{\mathrm{i}}|^2.
\]
Mit dem Kosinussatz:
\[
    q^2
    =
    p^2+p^2-2p^2\cos\theta
    =
    2p^2(1-\cos\theta).
\]
Mit
\[
    1-\cos\theta = 2\sin^2(\theta/2)
\]
folgt:
\[
    q^2 = 4p^2\sin^2(\theta/2).
\]
Damit:
\[
    q=2p\sin(\theta/2).
\]
\end{derivationbox}

\section{Strukturauflösung}

Die Auflösung eines Streuexperiments hängt von der Wellenlänge beziehungsweise vom
Impulsübertrag ab.

\begin{definitionbox}
Die \emph{Strukturauflösung} ist die kleinste Längenskala, die ein Experiment noch
unterscheiden kann.
\end{definitionbox}

\begin{formulabox}
Die de-Broglie-Wellenlänge eines Teilchens ist
\[
    \lambda = \frac{h}{p}.
\]
Für Streuprozesse ist die relevante Auflösung grob
\[
    \Delta x \sim \frac{\hbar}{q}.
\]
\end{formulabox}

\begin{conceptbox}
Je größer der Impulsübertrag $q$ ist, desto kleinere Strukturen können aufgelöst werden:
\[
    q \text{ groß}
    \quad\Rightarrow\quad
    \Delta x \text{ klein}.
\]

Das ist der Grund, warum hochenergetische Streuexperimente zur Untersuchung kleiner
Strukturen verwendet werden.
\end{conceptbox}

\begin{examtipbox}
Merksatz:
\[
    \text{hohe Energie}
    \Rightarrow
    \text{großer Impuls}
    \Rightarrow
    \text{großer Impulsübertrag möglich}
    \Rightarrow
    \text{kleine Strukturen auflösbar}.
\]
\end{examtipbox}

\section{Mottstreuung}

Die Rutherfordform ist klassisch und behandelt die streuenden Teilchen ohne Spin.
Für Elektronenstreuung an Kernen muss man relativistische Effekte und den Spin des
Elektrons berücksichtigen.

\begin{definitionbox}
Die \emph{Mottstreuung} ist die relativistische Streuung von Elektronen an einem
Coulombfeld unter Berücksichtigung des Elektronenspins.
\end{definitionbox}

\begin{conceptbox}
Mottstreuung ist besonders wichtig für die Bestimmung von Kernladungsverteilungen,
weil Elektronen punktförmig sind und elektromagnetisch wechselwirken.

Elektronen spüren hauptsächlich die elektrische Ladungsverteilung des Kerns.
Daher eignen sie sich als Sonden für den elektromagnetischen Radius.
\end{conceptbox}

\begin{formulabox}
Für ein punktförmiges spinloses Coulombzentrum ergibt sich gegenüber der Rutherfordform
eine relativistische/spinabhängige Korrektur.
Qualitativ kann man schreiben:
\[
    \left(\frac{\dd\sigma}{\dd\Omega}\right)_{\mathrm{Mott}}
    =
    \left(\frac{\dd\sigma}{\dd\Omega}\right)_{\mathrm{Ruth}}
    \left[
        1-\beta^2\sin^2\left(\frac{\theta}{2}\right)
    \right].
\]
Dabei ist
\[
    \beta=\frac{v}{c}.
\]
\end{formulabox}

\begin{notebox}
Je nach Konvention und Näherung erscheinen Mottformeln in leicht unterschiedlicher
Schreibweise.
Für dieses Kapitel ist die wichtigste Aussage:
Mottstreuung ist die passendere Referenz für Elektronenstreuung als die rein klassische
Rutherfordform.
\end{notebox}

\begin{examtipbox}
Rutherfordstreuung:
\begin{center}
klassische Coulombstreuung geladener Teilchen.
\end{center}

Mottstreuung:
\begin{center}
relativistische Coulombstreuung von Elektronen mit Spin.
\end{center}
\end{examtipbox}

\section{Formfaktor}

Wenn das Target nicht punktförmig ist, ändert sich die Streuverteilung.
Diese Änderung wird durch den Formfaktor beschrieben.

\begin{definitionbox}
Der \emph{Formfaktor} $F(q)$ beschreibt, wie die endliche räumliche Ausdehnung eines
Targets den Streuquerschnitt verändert.
\end{definitionbox}

\begin{formulabox}
Für ein ausgedehntes Target gilt häufig:
\[
    \left(\frac{\dd\sigma}{\dd\Omega}\right)_{\mathrm{real}}
    =
    \left(\frac{\dd\sigma}{\dd\Omega}\right)_{\mathrm{punkt}}
    |F(q)|^2.
\]
\end{formulabox}

\begin{conceptbox}
Der Punktquerschnitt beschreibt die Streuung an einem punktförmigen Objekt.

Der Faktor
\[
    |F(q)|^2
\]
enthält die Information über die innere räumliche Struktur.

Wenn das Target punktförmig ist, gilt:
\[
    F(q)=1.
\]
Dann ist
\[
    \left(\frac{\dd\sigma}{\dd\Omega}\right)_{\mathrm{real}}
    =
    \left(\frac{\dd\sigma}{\dd\Omega}\right)_{\mathrm{punkt}}.
\]
\end{conceptbox}

\section{Formfaktor und Ladungsverteilung}

Der Formfaktor ist eng mit der Ladungsverteilung verbunden.

\begin{definitionbox}
Die \emph{Ladungsverteilung} $\rho(\vec r)$ beschreibt, wie die elektrische Ladung
räumlich im Target verteilt ist.
\end{definitionbox}

\begin{formulabox}
Der Formfaktor ist im Wesentlichen die Fouriertransformierte der Ladungsverteilung:
\[
    F(\vec q)
    =
    \frac{1}{Q}
    \int \rho(\vec r)
    \exp\left(
        \frac{i\vec q\cdot\vec r}{\hbar}
    \right)
    \dd^3 r.
\]
Dabei ist
\[
    Q=\int \rho(\vec r)\,\dd^3 r
\]
die Gesamtladung.
\end{formulabox}

\begin{conceptbox}
Diese Formel bedeutet:
Streuung misst nicht direkt ein Bild des Kerns im Ortsraum.
Streuung misst zunächst Informationen im Impulsraum.

Aus der $q$-Abhängigkeit des Formfaktors kann man auf die räumliche Ladungsverteilung
zurückschließen.
\end{conceptbox}

\begin{notebox}
In den Vorlesungsabbildungen werden verschiedene Ladungsverteilungen mit ihren
Formfaktoren verglichen:
\begin{itemize}
    \item punktförmige Ladungsverteilung,
    \item exponentielle Verteilung,
    \item Gaußverteilung,
    \item homogene Kugel,
    \item Kugel mit diffuser Randzone.
\end{itemize}
Eine punktförmige Ladungsverteilung hat einen konstanten Formfaktor.
Ausgedehnte Ladungsverteilungen führen zu einem $q$-abhängigen Formfaktor.
\end{notebox}

\section{Elektromagnetischer Radius aus dem Formfaktor}

Bei kleinen Impulsüberträgen enthält der Formfaktor Information über den mittleren
quadratischen Radius.

\begin{definitionbox}
Der \emph{elektromagnetische Radius} ist ein Maß für die räumliche Ausdehnung der
elektrischen Ladungsverteilung.
Er wird häufig über den mittleren quadratischen Radius beschrieben:
\[
    r_{\mathrm{em}}
    =
    \sqrt{\langle r^2\rangle}.
\]
\end{definitionbox}

\begin{formulabox}
Für kleine Impulsüberträge gilt näherungsweise:
\[
    F(q)
    \approx
    1
    -
    \frac{q^2\langle r^2\rangle}{6\hbar^2}
    +\ldots
\]
Daraus kann man den mittleren quadratischen Radius bestimmen.
\end{formulabox}

\begin{conceptbox}
Die Steigung des Formfaktors bei kleinen $q$ liefert den Radius.

Ein punktförmiges Teilchen hat keinen ausgedehnten Ladungsradius und daher einen
konstanten Formfaktor.

Ein ausgedehntes Teilchen oder ein Kern hat einen Formfaktor, der mit $q$ variiert.
\end{conceptbox}

\begin{examtipbox}
Typische Prüfungsfrage:
\begin{center}
\emph{Was bedeutet ein konstanter Formfaktor?}
\end{center}

Antwort:
Ein konstanter Formfaktor ist ein Hinweis auf punktförmige Struktur auf der untersuchten
Längenskala.
\end{examtipbox}

\section{Beispiele für Formfaktoren}

\subsection{Punktförmiges Teilchen}

\begin{conceptbox}
Für ein punktförmiges Teilchen ist die Ladungsverteilung idealisiert:
\[
    \rho(\vec r)\propto \delta(\vec r).
\]
Dann ist der Formfaktor konstant:
\[
    F(q)=1.
\]
\end{conceptbox}

\subsection{Ausgedehnter Kern}

\begin{conceptbox}
Für einen ausgedehnten Kern ist die Ladung über ein endliches Volumen verteilt.

Dann interferieren die Streuamplituden von verschiedenen Raumpunkten.
Das führt zu einer $q$-abhängigen Modifikation der Streuung:
\[
    \left(\frac{\dd\sigma}{\dd\Omega}\right)_{\mathrm{real}}
    =
    \left(\frac{\dd\sigma}{\dd\Omega}\right)_{\mathrm{punkt}}
    |F(q)|^2.
\]
\end{conceptbox}

\subsection{Homogene Kugel}

\begin{notebox}
Für eine homogene Kugel zeigt der Formfaktor Oszillationen.
Diese Oszillationen hängen mit Interferenzeffekten zusammen.

Minima im Formfaktor entsprechen Impulsüberträgen, bei denen die Streuamplituden
aus verschiedenen Bereichen der Kugel destruktiv interferieren.
\end{notebox}

\section{Experimentelle Bedeutung}

\begin{conceptbox}
Streuexperimente liefern Informationen über:
\begin{itemize}
    \item Kernradien,
    \item Ladungsverteilungen,
    \item Formfaktoren,
    \item innere Struktur von Nukleonen,
    \item mögliche Substruktur von Teilchen.
\end{itemize}
\end{conceptbox}

\begin{notebox}
In den Vorlesungsfolien zur Proton-Substruktur wird gezeigt, dass bei hoher
Schwerpunktsenergie beziehungsweise großem Impulsübertrag der Formfaktor Hinweise
auf punktförmige Bestandteile liefert.

Das ist die experimentelle Grundlage für die Aussage, dass Protonen eine innere
Substruktur besitzen.
\end{notebox}

\section{Elastische und inelastische Streuung}

\begin{definitionbox}
Bei \emph{elastischer Streuung} bleiben die inneren Zustände der Teilchen unverändert.
Es ändern sich nur Impulse und Bewegungsrichtungen.

Bei \emph{inelastischer Streuung} wird zusätzlich innere Energie übertragen.
Das Target kann angeregt oder aufgebrochen werden.
\end{definitionbox}

\begin{conceptbox}
Elastische Streuung eignet sich besonders zur Untersuchung von Formfaktoren und
Ladungsverteilungen.

Inelastische Streuung kann tiefere Substruktur sichtbar machen.
Ein Beispiel ist die inelastische Elektron-Proton-Streuung, mit der die Substruktur
des Protons nachgewiesen wurde.
\end{conceptbox}

\section{Mini-Zusammenfassung}

\begin{notebox}
Die wichtigsten Punkte dieses Kapitels sind:
\begin{itemize}
    \item Streuung ist eine Methode zur Untersuchung kleiner Strukturen.
    \item Der differentielle Wirkungsquerschnitt beschreibt die Winkelverteilung.
    \item Rutherfordstreuung beschreibt klassische Coulombstreuung.
    \item Die Rutherfordform ist stark vorwärtsgerichtet:
    \[
        \frac{\dd\sigma}{\dd\Omega}
        \propto
        \frac{1}{\sin^4(\theta/2)}.
    \]
    \item Der Impulsübertrag ist bei elastischer Streuung
    \[
        q=2p\sin(\theta/2).
    \]
    \item Große Impulsüberträge erlauben kleine Strukturauflösung:
    \[
        \Delta x\sim \frac{\hbar}{q}.
    \]
    \item Mottstreuung ist die relativistische Elektronenstreuung mit Spin.
    \item Für ausgedehnte Targets gilt:
    \[
        \left(\frac{\dd\sigma}{\dd\Omega}\right)_{\mathrm{real}}
        =
        \left(\frac{\dd\sigma}{\dd\Omega}\right)_{\mathrm{punkt}}
        |F(q)|^2.
    \]
    \item Der Formfaktor enthält Information über die Ladungsverteilung.
    \item Ein konstanter Formfaktor bedeutet punktförmige Struktur auf der untersuchten Skala.
\end{itemize}
\end{notebox}

\section{Prüfungscheck}

\begin{examtipbox}
Nach diesem Kapitel solltest du folgende Fragen beantworten können:
\begin{enumerate}
    \item Was misst ein Streuexperiment?
    \item Was ist ein differentieller Wirkungsquerschnitt?
    \item Warum ist Rutherfordstreuung stark vorwärtsgerichtet?
    \item Welche Annahmen stecken in der Rutherfordstreuung?
    \item Was bedeutet eine Abweichung von der Rutherfordform?
    \item Was ist der Rutherfordradius?
    \item Was ist der Impulsübertrag?
    \item Wie lautet $q$ bei elastischer Streuung?
    \item Warum verbessert großer Impulsübertrag die Strukturauflösung?
    \item Was ist Mottstreuung?
    \item Was ist ein Formfaktor?
    \item Warum gilt für ausgedehnte Targets ein Faktor $|F(q)|^2$?
    \item Was ist der Zusammenhang zwischen Formfaktor und Ladungsverteilung?
    \item Was bedeutet ein konstanter Formfaktor?
    \item Wie erhält man aus dem Formfaktor den elektromagnetischen Radius?
    \item Was ist der Unterschied zwischen elastischer und inelastischer Streuung?
\end{enumerate}
\end{examtipbox}